\textbf{Phần 1. Trắc nghiệm lựa chọn}
\setcounterpageref{ex}{0}

\begin{ex} %[2H3V17]
Để có đúng $7\text{ kg}$ gạo từ bao gạo to, người bán hàng dùng cân hai đĩa với quả cân loại nào?
\choice
{$3$ quả $2\text{ kg}$}
{$1$ quả $5\text{ kg}$ và $1$ quả $1\text{ kg}$}
{\True $1$ quả $5\text{ kg}$ và $1$ quả $2\text{ kg}$}
{$2$ quả $3\text{ kg}$}
\loigiai{
$5 + 2 = 7\text{ (kg)}$.
\textbf{Đáp án của câu hỏi là C}
}
\end{ex}

\begin{ex} %[2H3V17]
Một bình đầy nước có thể rót đầy $4$ ca nước loại $1\text{ L}$ và $1$ ca nước loại $2\text{ L}$. Bình nước đó chứa được bao nhiêu lít nước?
\choice
{$5\text{ L}$}
{\True $6\text{ L}$}
{$7\text{ L}$}
{$8\text{ L}$}
\loigiai{
$4 \times 1 + 2 = 6\text{ (L)}$.
\textbf{Đáp án của câu hỏi là B}
}
\end{ex}

\begin{ex} %[2H3V17]
Gia đình Mai mỗi ngày ăn hết $1\text{ kg}$ gạo. Hỏi trong $2$ tuần gia đình Mai ăn hết bao nhiêu ki-lô-gam gạo?
\choice
{$12\text{ kg}$}
{$10\text{ kg}$}
{$16\text{ kg}$}
{\True $14\text{ kg}$}
\loigiai{
$2$ tuần $= 14$ ngày. Gia đình Mai ăn hết: $14 \times 1 = 14\text{ (kg)}$.
\textbf{Đáp án của câu hỏi là D}
}
\end{ex}

\begin{ex} %[2H3V17]
Có $1$ can $3\text{ L}$ và $1$ can $5\text{ L}$. Làm thế nào để lấy được đúng $2\text{ L}$ nước từ thùng nước lớn?
\choice
{\True Múc đầy can $5\text{ L}$, rót sang cho đầy can $3\text{ L}$, lượng nước còn lại trong can $5\text{ L}$ là $2\text{ L}$}
{Múc $2$ lần can $3\text{ L}$}
{Múc đầy can $3\text{ L}$ đổ vào can $5\text{ L}$}
{Múc nửa can $5\text{ L}$}
\loigiai{
Lượng nước còn lại trong can 5L là: $5 - 3 = 2\text{ (L)}$.
\textbf{Đáp án của câu hỏi là A}
}
\end{ex}

\begin{ex} %[2H3V17]
Bác Ba có $24\text{ kg}$ xoài, bác đã bán đi $16\text{ kg}$. Số xoài còn lại bác chia đều vào $2$ rổ. Mỗi rổ có bao nhiêu ki-lô-gam xoài?
\choice
{$3\text{ kg}$}
{$5\text{ kg}$}
{\True $4\text{ kg}$}
{$6\text{ kg}$}
\loigiai{
Số xoài còn lại: $24 - 16 = 8\text{ (kg)}$. Mỗi rổ có: $8 : 2 = 4\text{ (kg)}$.
\textbf{Đáp án của câu hỏi là C}
}
\end{ex}

\begin{ex} %[2H3V17]
Một thùng sơn nặng $18\text{ kg}$. Sơn xong một bức tường, người ta thấy thùng sơn còn lại $9\text{ kg}$. Khối lượng sơn đã dùng là
\choice
{$8\text{ kg}$}
{\True $9\text{ kg}$}
{$10\text{ kg}$}
{$7\text{ kg}$}
\loigiai{
$18 - 9 = 9\text{ (kg)}$.
\textbf{Đáp án của câu hỏi là B}
}
\end{ex}

\begin{ex} %[2H3V17]
Mẹ đong nước mắm vào các chai nhỏ loại $2\text{ L}$. Có tất cả $16\text{ L}$ nước mắm. Mẹ đong được bao nhiêu chai?
\choice
{$7$ chai}
{$6$ chai}
{$9$ chai}
{\True $8$ chai}
\loigiai{
$16 : 2 = 8$ (chai).
\textbf{Đáp án của câu hỏi là D}
}
\end{ex}

\begin{ex} %[2H3V17]
Con ngỗng nặng $5\text{ kg}$, con vịt nhẹ hơn con ngỗng $2\text{ kg}$, con gà nặng bằng con vịt. Hỏi cả ba con nặng bao nhiêu ki-lô-gam?
\choice
{\True $11\text{ kg}$}
{$10\text{ kg}$}
{$12\text{ kg}$}
{$9\text{ kg}$}
\loigiai{
Con vịt nặng: $5 - 2 = 3\text{ (kg)}$. Con gà nặng $3\text{ kg}$. Cả ba con nặng: $5 + 3 + 3 = 11\text{ (kg)}$.
\textbf{Đáp án của câu hỏi là A}
}
\end{ex}

\begin{ex} %[2H3V17]
Một bình trà sữa có $15\text{ L}$. Bán đi một số lít thì còn lại $8\text{ L}$. Đã bán bao nhiêu lít trà sữa?
\choice
{$6\text{ L}$}
{$8\text{ L}$}
{\True $7\text{ L}$}
{$9\text{ L}$}
\loigiai{
$15 - 8 = 7\text{ (L)}$.
\textbf{Đáp án của câu hỏi là C}
}
\end{ex}

\begin{ex} %[2H3V17]
Túi thứ nhất đựng $12\text{ kg}$ khoai tây, túi thứ hai đựng nhiều hơn túi thứ nhất $5\text{ kg}$. Cả hai túi đựng bao nhiêu ki-lô-gam khoai tây?
\choice
{$28\text{ kg}$}
{\True $29\text{ kg}$}
{$30\text{ kg}$}
{$27\text{ kg}$}
\loigiai{
Túi 2: $12 + 5 = 17\text{ (kg)}$. Cả hai túi: $12 + 17 = 29\text{ (kg)}$.
\textbf{Đáp án của câu hỏi là B}
}
\end{ex}

\begin{ex} %[2H3V17]
Điền số thích hợp: $35\text{ L} - 17\text{ L} = 9\text{ L} + \dots\text{ L}$.
\choice
{$8$}
{$10$}
{$7$}
{\True $9$}
\loigiai{
Vế trái: $35 - 17 = 18\text{ L}$. Vế phải: $18 - 9 = 9$. Vậy điền 9.
\textbf{Đáp án của câu hỏi là D}
}
\end{ex}

\begin{ex} %[2H3V17]
Muốn có $10\text{ L}$ nước, người ta có thể dùng những chiếc can loại nào dưới đây?
\choice
{\True $2$ chiếc can $5\text{ L}$}
{$3$ chiếc can $3\text{ L}$}
{$1$ can $5\text{ L}$ và $2$ can $2\text{ L}$}
{$4$ can $2\text{ L}$}
\loigiai{
$2 \times 5 = 10\text{ (L)}$.
\textbf{Đáp án của câu hỏi là A}
}
\end{ex}

\textbf{Phần 2. Trắc nghiệm đúng sai}
\setcounterpageref{ex}{0}

\begin{ex} %[2H3V17]
Xét tính đúng sai của các tình huống thực hành cân đong:
\choiceTF
{\True Đĩa cân bên trái có gói đường và quả cân 2 kg, đĩa bên phải có quả cân 5 kg. Cân thăng bằng thì gói đường nặng 3 kg.}
{Chỉ dùng can 2 L và can 5 L thì không thể đong được 7 L nước.}
{\True Tổng khối lượng của một vật và đồ đựng luôn lớn hơn khối lượng của riêng vật đó.}
{Một lít nước nặng hơn 100 kg bông gòn.}
\loigiai{
\begin{itemize}
    \item Ý a) Đúng: $5 - 2 = 3\text{ kg}$. Đúng.
    \item Ý b) Sai: $2 + 5 = 7\text{ L}$. Hoàn toàn đong được.
    \item Ý c) Đúng: Khối lượng cả bì = Khối lượng vật + Khối lượng đồ đựng.
    \item Ý d) Sai: 1 lít nước nặng khoảng 1 kg, nhẹ hơn rất nhiều so với 100 kg bông gòn.
\end{itemize}
\begin{itemize}
    \item a) \textbf{Đúng}.
    \item b) \textbf{Sai}.
    \item c) \textbf{Đúng}.
    \item d) \textbf{Sai}.
\end{itemize}
}
\end{ex}

\begin{ex} %[2H3V17]
Bác thợ mộc có hai thanh gỗ: thanh thứ nhất nặng 18 kg, thanh thứ hai nặng 14 kg:
\choiceTF
{\True Thanh thứ nhất nặng hơn thanh thứ hai 4 kg.}
{\True Cả hai thanh gỗ nặng 32 kg.}
{Nếu cưa bớt ở thanh thứ nhất 4 kg thì hai thanh có khối lượng bằng nhau và bằng 18 kg.}
{\True Nếu ghép hai thanh lại rồi cưa đôi thì mỗi nửa nặng 16 kg.}
\loigiai{
\begin{itemize}
    \item Ý a) Đúng: $18 - 14 = 4\text{ kg}$.
    \item Ý b) Đúng: $18 + 14 = 32\text{ kg}$.
    \item Ý c) Sai: Khi đó cả hai thanh đều nặng 14 kg (không phải 18 kg).
    \item Ý d) Đúng: $32 : 2 = 16\text{ kg}$. Đúng.
\end{itemize}
\begin{itemize}
    \item a) \textbf{Đúng}.
    \item b) \textbf{Đúng}.
    \item c) \textbf{Sai}.
    \item d) \textbf{Đúng}.
\end{itemize}
}
\end{ex}

\textbf{Phần 3. Trả lời ngắn}
\setcounterpageref{ex}{0}

\begin{ex} %[2H3V17]
Một người thợ rót mật ong từ can 10 lít vào các chai nhỏ loại 1 lít. Rót được 6 chai thì can còn lại bao nhiêu lít mật ong?
\shortans{4}
\loigiai{
\textbf{Đáp số:} $4$\\
$10 - 6 = 4\text{ L}$.
}
\end{ex}

\begin{ex} %[2H3V17]
Một con cá quả nặng 2 kg, con cá chép nặng hơn con cá quả 1 kg. Hỏi cả hai con cá nặng bao nhiêu ki-lô-gam?
\shortans{5}
\loigiai{
\textbf{Đáp số:} $5$\\
Cá chép: $2 + 1 = 3\text{ kg}$. Tổng: $2 + 3 = 5\text{ kg}$.
}
\end{ex}

\begin{ex} %[2H3V17]
Có 3 bao ngô, mỗi bao nặng 9 kg. Hỏi cả 3 bao ngô nặng tất cả bao nhiêu ki-lô-gam?
\shortans{27}
\loigiai{
\textbf{Đáp số:} $27$\\
$9 \times 3 = 27\text{ kg}$.
}
\end{ex}

\begin{ex} %[2H3V17]
Một thùng có 25 lít nước mắm, người ta đã bán đi 16 lít nước mắm. Hỏi trong thùng còn lại bao nhiêu lít nước mắm?
\shortans{9}
\loigiai{
\textbf{Đáp số:} $9$\\
$25 - 16 = 9\text{ L}$.
}
\end{ex}

\begin{ex} %[2H3V17]
Một túi kẹo nặng 1 kg, người ta đóng 10 túi như thế vào một hộp. Hỏi 10 túi kẹo nặng bao nhiêu ki-lô-gam?
\shortans{10}
\loigiai{
\textbf{Đáp số:} $10$\\
$1 \times 10 = 10\text{ kg}$.
}
\end{ex}

\begin{ex} %[2H3V17]
Một can 20 lít đựng đầy xăng. Sau khi rót vào xe máy 6 lít và xe ô tô 9 lít thì trong can còn lại bao nhiêu lít xăng?
\shortans{5}
\loigiai{
\textbf{Đáp số:} $5$\\
$20 - 6 - 9 = 5\text{ L}$.
}
\end{ex}
